\section{Introduction}
\label{sec:intro}

National federations choose tournament rosters from a pool of professional players whose club-level performance is well measured, but whose performance together has rarely or never been observed. Selection committees combine scouting, statistics and judgment, and the reasoning behind a final roster is usually not written down in a form that can be checked. This paper does not claim that a statistical model can replace that judgment. The quantity a committee would most like to know, the best possible roster for a short tournament, is counterfactual: it is never observed, so no method can be validated by comparing its output to it. What a model can do is make one specific set of assumptions explicit, apply them consistently to every candidate roster, and report where and by how much its preferences differ from the roster that was actually chosen. We call this an audit of the selection, and we design every part of the method and its evaluation around that use.

The method has three parts. A graph neural network, trained on regular-season games, predicts the outcome of a game from the players on each side and from how much time each pair of players spent on the floor or ice together. A deployment model turns a candidate roster into those shared-time matrices by mixing over the units a coach can actually play (five-player lineups in basketball, forward lines and defensive pairs in hockey). Finally, an adversarial evaluation treats deployment as a leader-follower game in the sense of \citet{conitzer2006}: the candidate roster commits to a deployment, and each reference opponent chooses its own deployment in response. A roster's score is its predicted win probability under this game, averaged over a fixed set of opponents. A search procedure explores the space of rosters that satisfy positional requirements, and a replacement-level knockout attributes the score to individual players.

We apply the method to two tournaments: men's basketball at the 2024 Paris Olympics (United States and Canada) and the 2025 4 Nations Face-Off in hockey (Canada). For each case we report how the model scores the roster that was actually selected, how far each selected player is from being chosen or dropped by the model, how the model's preferred roster compares with simple baselines under the same constraints, and how stable these conclusions are across model refits, search seeds and modeling choices. The main contributions are (i) a deployment model that keeps adversarial evaluation inside the space of realizable lineups, (ii) an evaluation protocol for a counterfactual selection problem built from out-of-sample prediction checks, baseline comparisons, null distributions and sensitivity analysis rather than agreement with expert choices, and (iii) evidence on which parts of the pipeline carry information and which do not.

Two choices in this framework involve domain judgment and we keep both explicit. Positional quotas restrict the search to rosters that contain enough players at each position, and some analyses lock a small number of players into the roster. Section~\ref{sec:problem} defines both, Section~\ref{sec:cases} reports every analysis with and without locks, and Section~\ref{sec:sensitivity} reports the effect of removing the quotas. The method is useful to the extent that it changes or sharpens a selection discussion, and we return to the limits of that use in Section~\ref{sec:discussion}.
